\cotefile{43}
\scanpage{32}
% archive sheet number: 31
% s1, lines 1--3
{\unsure{Utilisant}} une technique de descente analogue, on \deleted{{\unsure{et}}} définit {\unsure{par descente}} un homomorphisme
\[
H^1(X,G)\xrightarrow{d}\operatorname{Hom}_{Y\text{-}\mathrm{grp}}(G_m,G)
\qquad (\uline{\text{{\unsure{hom degré|hom. de gr.}}}})
\]
% margin note slanted; two-ended reference line points to the preceding display
\margin{{\unsure{la définition dans le cas général à voir au cas}} $Y=\operatorname{Spec}(k)$, $k$ {\unsure{alg. clos}}, $G=G_m$ {\unsure{\ldots}}}
\diagram{\begin{tikzcd}[ampersand replacement=\&]{}\arrow[r,leftrightarrow]\&{}\end{tikzcd}}
et on trouve que {\unsure{la suite}}
\[
\boxed{0\longrightarrow \operatorname{Ext}^1({\unsure{\mathscr{P}|\Phi}}_{X/{\unsure{Y|k}}},G)
\longrightarrow H^1({\unsure{X_S|X,S|X}},G)
\longrightarrow\operatorname{Hom}_{Y\text{-}\mathrm{grp}}(G_m,G)}
\]
est exacte. {\unsure{Il est même plausible que}} cette suite s'insère dans la suite exacte prolongée
\[
\boxed{\longrightarrow\operatorname{Ext}^2(\mathscr{P}_{X/Y},G)
\longrightarrow H^2(X,G)\longrightarrow\operatorname{Ext}^1_{Y\text{-}\mathrm{grp}}(G_m,G)\longrightarrow\cdots}
\]
où {\unsure{on fera paraître|a faire paraître}} l'hom
\[
\operatorname{Ext}^i(G_m,G)\longrightarrow\operatorname{Ext}^{i+2}(\mathscr{P}_{X/Y},G)
\]
est la multiplication par un élément fondamental
\[
{\unsure{\zeta|\xi}}\in\operatorname{Ext}^2_{Y\text{-}\mathrm{grp}}(\mathscr{P}_{X/Y},G_m)
\]
dont la définition {\unsure{directe}}; la formule {\unsure{serait}} la suivante [{\unsure{valable}} s'il existe un $S$ {\unsure{étale}} tq $S\to Y$ soit surjectif,
